\section{Institutional Background}

Brazil's federal SME statute (Complementary Law 123/2006, amended 2014) defines micro- and small enterprises by annual gross revenue---R\$360{,}000 (US\$103{,}000) and R\$4{,}800{,}000 (US\$1.37 million) ceilings, respectively---and mandates exclusive SME tenders for items valued at R\$80{,}000 (US\$22{,}900) or less. Public buyer units (PBUs)---federal, state, and municipal administrative entities subject to procurement rules---can avoid the SME-only default only by justifying the departure through a formal report to the audit courts, a costly and scrutinized process. The 2014 amendment converted what had been an option (``may'') into a binding default (``must'') and imposed these \emph{opt-out costs}; state-level adherence to SME-only tendering for eligible items rose from roughly 13\% to nearly 70\% in the subsequent years.

Since 2005, all public procurement in the state of Sao Paulo passes through Bolsa Eletr\^onica de Compras (BEC), a centralized electronic reverse-auction platform in which registered suppliers submit sealed bids (phase 1) followed by an iterative electronic auction (phases 2 onward) in which each participant may progressively lower its offer until the auction closes. BEC handled approximately R\$13 billion (US\$3.7 billion) of procurement in 2019 alone and has cumulatively moved more than R\$105 billion (US\$30 billion) since implementation, spanning 860{,}000 purchase offers and 4.8 million item-level transactions. BEC organizes its catalog into three tiers---group, class, and item---so that ``medical, dental, and hospital equipment and supplies'' form Group 65, within which class 6531 identifies medications subject to price ceilings set by the C\^amara de Regula\c{c}\~ao do Mercado de Medicamentos (CMED).

Group 65 was a singular exception to the 2014 rule. Between August 2014 and February 2018, the state government and the Tribunal de Contas do Estado (TCE-SP) operated under a joint interpretation that medical supplies constituted strategic items requiring open competition; no bid for Group 65 items was restricted to SMEs during this period. In November 2017, the Procuradoria-Geral do Estado (PGE-SP, a separate control agency) issued a legal opinion reversing this exception on grounds of the isonomy principle---the constitutional requirement that comparable cases be treated alike---arguing that no substantive feature of medical supplies justified procedural distinctions from other product groups under state procurement law. Effective March 2018, opt-out costs extended to Group 65; adherence to SME-only tendering within the group jumped to 43\% on average, below the rate for other groups largely because class-6531 pharmaceuticals are oligopolistic and often fail the requirement of at least three eligible SME suppliers.

The critical feature of this institutional timeline for identification is that the policy change was triggered by a legalistic reinterpretation of the isonomy principle by PGE-SP rather than by deteriorating procurement outcomes for Group 65. Public documentation reveals no evidence that the decision responded to price levels, competition patterns, or supplier complaints specific to Group 65 in 2017. The timing of implementation was driven by bureaucratic processes within PGE-SP and its communication to PBUs, making the change plausibly exogenous to the outcomes studied. These features provide the institutional basis for the identifying assumption that, absent the March 2018 reinterpretation, Group 65 would have followed parallel outcome trajectories with the 76 never-treated product groups in the post-period.
